\documentclass[11pt]{article}
\usepackage[letterpaper,margin=0.72in]{geometry}
\usepackage{amsmath,amssymb,amsfonts,bm,booktabs,array,enumitem}
\usepackage[colorlinks=true,linkcolor=blue,citecolor=blue,urlcolor=blue]{hyperref}

\setlength{\parindent}{0pt}
\setlength{\parskip}{4pt}
\setlist[itemize]{leftmargin=1.2em,itemsep=1pt,topsep=2pt}
\setlist[enumerate]{leftmargin=1.4em,itemsep=1pt,topsep=2pt}
\newcommand{\innerT}[2]{\langle #1,#2\rangle_T}
\newcommand{\innerP}[2]{\langle #1,#2\rangle_P}
\newcommand{\E}{\mathbb E}
\newcommand{\N}{\mathcal N}

\pagestyle{empty}
\begin{document}

\begin{center}
{\Large\bf Profile Sieve EL with Nuisance Breaks}\\[-1pt]
{\large A Geometry-First Proof Map}\\[-1pt]
\end{center}

\vspace{-4pt}
\textbf{One-sentence thesis.}
Break-adaptive projection is the decontamination step that turns false
predictability from structural breaks into valid empirical-likelihood inference
for \(\beta\).

\section*{1. Gap: structural breaks can look like predictability}

The predictive model is
\[
    Y_t=\beta_0X_{t-1}+\mu_{0,t}+U_t,\qquad
    \mu_{0,t}=m_j(W_{t-1})\quad\text{on regime }j .
\]
The difficulty is not only that \(\mu_{0,t}\) is nonparametric. It may also
break. If inference residualizes with a stable/no-break nuisance space, its
residual-maker \(M_{\rm st}\) generally leaves a nonzero nuisance residue:
\[
    M_{\rm st}\mu_0\neq0 .
\]
At the true value \(\beta_0\), the stable projected score is
\[
    S_T^{\rm st}(\beta_0)
    =
    T^{-1/2}u'M_{\rm st}w
    +
    T^{-1/2}\mu_0'M_{\rm st}w .
\]
The second term is omitted-break drift:
\[
    B_T
    =
    T^{-1/2}\mu_0'M_{\rm st}w
    =
    \sqrt T\,\innerT{M_{\rm st}\mu_0}{M_{\rm st}w}.
\]
Thus false predictability is a geometric angle problem. If the unremoved
nuisance residue \(M_{\rm st}\mu_0\) is not orthogonal to the residualized score
direction \(M_{\rm st}w\), the score can be large even when \(\beta=\beta_0\).

\textbf{Gap.} Existing stable-nuisance predictability procedures can confuse
nuisance structural breaks with predictive signal. The missing piece is a
method that ties break selection, nuisance removal, and EL inference together
so that break residue cannot enter the predictive score.

\textbf{New idea.} Treat a candidate break partition as a nuisance subspace:
\[
    \Lambda\longmapsto
    \N_{K,\Lambda,T}=\operatorname{col}(Q_{K,\Lambda})\subset L^2(P_T),
    \qquad
    M_{K,\Lambda}=I-\Pi_{K,\Lambda,T}.
\]
Break search is then subspace selection, not just date estimation.

\textbf{Why it matters.} Under the correct break-aware subspace,
\[
    M_0\mu_0\approx0,
\]
whereas under a stable nuisance subspace \(M_{\rm st}\mu_0\not\approx0\). The
geometry identifies exactly which component must be removed before testing
predictability.

\newpage
\section*{2. Geometry: break search and inference happen in normal space}

Let \(L^2(P_T)=(\mathbb R^T,\innerT{a}{b}=T^{-1}a'b)\). For a trial
\(\beta\), write \(y_\beta=y-\beta x\). Profiling over a block sieve nuisance
is orthogonal projection:
\[
    RSS_K(\beta,\Lambda)
    =
    T\operatorname{dist}_T^2(y-\beta x,\N_{K,\Lambda,T})
    =
    T\|M_{K,\Lambda}(y-\beta x)\|_T^2 .
\]
So the break estimator chooses the nuisance subspace closest to \(y_\beta\).
The projected score uses only normal-space components:
\[
    \Psi_T(\beta,\Lambda)
    =
    \innerT{M_{K,\Lambda}(y-\beta x)}
            {M_{K,\Lambda}w}.
\]
This is the decontaminated predictability moment.

\textbf{Population identification.} Let \(M_0^P\) project away the true
break-aware nuisance tangent space. Since \(\mu_0\) belongs to that nuisance
space,
\[
    M_0^P\mu_0=0 .
\]
The population moment satisfies
\[
    \Psi(\beta,\Lambda_0)
    =
    \innerP{M_0^P(Y-\beta X)}{M_0^Pw}
    =
    -(\beta-\beta_0)\Gamma,
    \qquad
    \Gamma=\innerP{M_0^PX}{M_0^Pw}\neq0.
\]
Thus \(\beta_0\) is identified by transversality: after quotienting out the
nuisance space, the residualized parametric direction is not orthogonal to the
chosen score direction.

\textbf{Good-event proof map.} Probability builds the random geometry. Once the
following good event holds, the proof is deterministic:
\[
    \mathcal G_T
    =
    G_T^{iso}\cap G_T^{approx}\cap G_T^{angle}
    \cap G_T^{local}\cap G_T^{select}\cap G_T^{convex}.
\]
\[
\begin{array}{p{0.22\linewidth}p{0.69\linewidth}}
\toprule
Component & Meaning and use\\
\midrule
\(G_T^{iso}\) & Gram and inner-product stability for projection and variance.\\
\(G_T^{approx}\) & \(M_0\mu_0\) is small, so nuisance residue is negligible.\\
\(G_T^{angle}\) & \(\Gamma\neq0\) and score variance is nondegenerate.\\
\(G_T^{local}\) & \(\widehat\Lambda\) replaces \(\Lambda_0\) in score and variance.\\
\(G_T^{select}\) & RSS localizes breaks and selects the order.\\
\(G_T^{convex}\) & EL convex hull and max-score conditions hold.\\
\bottomrule
\end{array}
\]
The full stochastic proof verifies \(\Pr(\mathcal G_T)\to1\) through Gram
laws, sieve approximation, RSS concentration, break replacement, CLT, variance
LLN, max-score, and convex-hull bounds.

\newpage
\section*{3. EL inference and the payoff}

For fixed \((\beta,\Lambda)\), define
\[
    g_t(\beta,\Lambda)
    =
    \{M_{K,\Lambda}(y-\beta x)\}_t\{M_{K,\Lambda}w\}_t .
\]
The scalar empirical likelihood is a convex simplex problem:
\[
    \max_{\pi\in\Delta_T}\sum_{t=1}^T\log(T\pi_t)
    \quad\text{s.t.}\quad
    \sum_{t=1}^T\pi_t g_t(\beta,\Lambda)=0.
\]
Its KKT form is
\[
    \pi_t(\lambda)=\frac{1}{T\{1+\lambda g_t\}},
    \qquad
    P_T\left\{\frac{g_t}{1+\lambda g_t}\right\}=0.
\]
On \(\mathcal G_T\), the logarithmic barrier has the local quadratic expansion
\[
    -2\log EL(\beta_0,\widehat\Lambda)
    =
    \frac{T\{\Psi_T(\beta_0,\widehat\Lambda)\}^2}
         {V_T(\widehat\Lambda)}
    +o_p(1).
\]
The break-aware projection bridge gives the oracle score reduction
\[
    \sqrt T\,\Psi_T(\beta_0,\widehat\Lambda)
    =
    T^{-1/2}u'M_0w+o_p(1).
\]
If
\[
    T^{-1/2}u'M_0w\Rightarrow N(0,V),\qquad
    V_T(\widehat\Lambda)\to_p V>0,
\]
then
\[
    -2\log EL(\beta_0,\widehat\Lambda)\Rightarrow\chi_1^2 .
\]

\textbf{Why this fills the gap.} The method turns unknown nuisance breaks into
a subspace-selection problem:
\[
    \text{break partition}
    \Rightarrow
    \text{nuisance subspace}
    \Rightarrow
    \text{orthogonal residual}
    \Rightarrow
    \text{projected score}
    \Rightarrow
    \text{EL KKT}
    \Rightarrow
    \chi_1^2 .
\]
The contaminated stable score is replaced by the oracle break-aware score. That
is why the geometry is central to the paper, not a notation choice.

\textbf{Efficiency add-on.} Under a conditional Gaussian location efficiency
calculation with variance \(\sigma^2(S,W)\), choose
\[
    w^*=\frac{X}{\sigma^2(S,W)}.
\]
Then
\[
    M_0^Pw^*
    =
    \frac{X-E(X\mid S,W)}{\sigma^2(S,W)},
\]
so the projected score equals the efficient score. In the homoskedastic case,
using \(w=X\) gives the same root.

\end{document}
